%-*-latex-*-
%-*-latex-*-
\newcommand\COURSE{ciss350}
\newcommand\ASSESSMENT{q10701}
\newcommand\ASSESSMENTTYPE{Quiz}
\newcommand\POINTS{\textwhite{xxx/xxx}}

\input{myquizpreamble}
\input{yliow}
\input{\COURSE}
%-*-latex-*-

%-*-latex-*-

%-*-latex-*-
\renewcommand\TITLE{\ASSESSMENTTYPE \ \ASSESSMENT}

\renewcommand\EMAIL{}
\renewcommand\AUTHOR{eagraham1@cougars.ccis.edu}

\textwidth=6in
\begin{document}
\topmatterthree

You are given the following binary tree node class:
\begin{Verbatim}[frame=single,fontsize=\footnotesize]
class BTNode
{
public:
    int key_;
    BTNode * parent_, * left_, * right_;
};
\end{Verbatim}
The following questions uses this class.

(You should create test cases to check that your functions work.)

%------------------------------------------------------------------------------
\nextq
Complete the function below that returns the pointer
pointing to the sibling node.
For instance of a node \verb!n1! has left child \verb!n2! and
right child \verb!n3!,
and \verb!p! points to node \verb!n2!, then \verb!sibling(p)! will return
the address of \verb!n3!.
If there is no sibling, \verb!NULL! is returned.
(For instance if \verb!p! points to the root node, \verb!NULL!
is returned.
Also, if the parent of \verb!*p! has only one child, then \verb!NULL! is
also returned.)
\\
\ANSWER
\begin{answercode}
BTNode * sibling(BTNode * p)
{
  if (p == NULL)
      return NULL;
  
  BTNode * q = p->parent_;

  if(q->left_ == p)
      return q->right_;
  if(q->right_ == p)
      return q->left_;
  else
      return NULL;
}
\end{answercode}

%------------------------------------------------------------------------------
\nextq
Complete the following function that returns the size of the
tree that \verb!p! points to (i.e., \verb!p! points to the root of this tree).
\\
\ANSWER
\begin{answercode}
int size(BTNode * p)
{
  if (p == NULL)
      return 0;

  BTNode * left = p->left_;
  BTNode * right = p->right_;

  return 1 + (size(left) + size(right));
}
\end{answercode}

%------------------------------------------------------------------------------
\nextq
Complete the following function that returns the maximum value of the
tree that \verb!p! points to (i.e., \verb!p! points to the root of this tree).
In other words, it returns the maximum of all the \verb!key_! values
of the nodes in the tree that \verb!p! points to.
\\
\ANSWER
\begin{answercode}
int max(BTNode * p)
{
  if(p == NULL)
      return 0;

  int leftMax = max(p->left_);
  int rightMax = max(p->right_);

  int m = p->key_;

  if(leftMax > m) m = leftMax;
  if(rightMax > m) m = rightMax;

  return m;
}
\end{answercode}

%------------------------------------------------------------------------------
%-*-latex-*-
\newpage
\input{instructions}
\end{document}

